% Working notes for MP2 (Introduction to Robotics, Group 6).
% Not the final report: collect facts, decisions and results here while working,
% then shorten them into the report.
\documentclass[10pt]{article}
\usepackage[margin=2cm]{geometry}
\usepackage{amsmath,amssymb}
\usepackage{booktabs}
\usepackage{graphicx}
\usepackage{xcolor}
\usepackage{hyperref}
\newcommand{\todo}[1]{\textcolor{red}{\textbf{TODO:} #1}}
\newcommand{\pkg}[1]{\texttt{#1}}

\title{MP2 -- Working Notes}
\author{Group 6}
\date{started 2026-11-10}

\begin{document}
\maketitle

% =====================================================================
\section{Environment}
\begin{itemize}
  \item ROS~2 Jazzy, Gazebo Harmonic (\pkg{gz sim}), Ubuntu 24.04 in WSL2 (WSLg for the GUI).
  \item Gazebo server and GUI start under WSL2 without errors (tested 2026-11-10).
  \item \pkg{TURTLEBOT3\_MODEL=waffle\_pi}, \pkg{ROS\_DOMAIN\_ID=30}.
  \item Real robot for Task~2: TurtleBot3 Waffle Pi or TurtleBot4 (our choice). Proposed: TurtleBot3 Waffle Pi
        (map, AMCL parameters and experience from MP1; same model as in the simulation). \todo{confirm}
  \item Deadline: end of December 2026.
\end{itemize}

% =====================================================================
\section{Task 1a -- Global path planning and path following in simulation}

\subsection{Goal}
Run Nav2 in Gazebo, set the initial pose and the goal in RViz, and explain how the
components work together: global planner (path), local planner/controller (velocity
commands), costmaps, \pkg{amcl}, \pkg{map\_server}. Then vary parameters and compare
the resulting paths with quantitative metrics.

\subsection{Simulation setup}
\textbf{Chosen: ROBOTIS simulation} (\pkg{ros-jazzy-turtlebot3-simulations} 2.3.7), as in the
tutorial of the assignment: \pkg{turtlebot3\_gazebo} (world) + \pkg{turtlebot3\_navigation2/navigation2.launch.py}
(Nav2 with \pkg{param/waffle\_pi.yaml}).
Reasons: follows the assignment, robot model Waffle Pi = real robot, DWB (DWA) as default
controller, and the same parameter file is used for the real robot in Task~2.

For the experiments: own launch file \pkg{mp2/launch/sim.launch.py} with the same world, robot and Nav2 bringup
(\pkg{bringup\_launch.py} with \pkg{mp2/config/nav2\_params.yaml}), headless by default, robot model without
camera and with ground-truth pose publisher (see part 2).

Alternative (not used): \pkg{nav2\_bringup/tb3\_simulation\_launch.py} starts Gazebo, Nav2 and RViz in one launch,
world \pkg{tb3\_sandbox}, model \emph{waffle}, parameters \pkg{nav2\_bringup/params/nav2\_params.yaml} (MPPI controller).

\subsection{Default configuration (as installed)}
\begin{tabular}{lll}
\toprule
Component & \pkg{waffle\_pi.yaml} (ROBOTIS) & \pkg{nav2\_params.yaml} (Nav2) \\
\midrule
Global planner & NavFn, \pkg{use\_astar: false} (Dijkstra) & NavFn, \pkg{use\_astar: false} (Dijkstra) \\
Controller (local planner) & DWB (\pkg{dwb\_core::DWBLocalPlanner}) & MPPI \\
\bottomrule
\end{tabular}

\paragraph{TEB.} The Timed Elastic Band controller is not released for ROS~2 Jazzy.
It is explained in theory only; in practice DWB is compared with MPPI and/or
Regulated Pure Pursuit (both installed). \textbf{TEB:} optimizes the trajectory as a sequence of timed poses
(``elastic band'') with respect to travel time, obstacle distance and kinematic limits.

\subsection{Overview: Gazebo vs.\ Nav2 vs.\ RViz}
\begin{itemize}
  \item \textbf{Gazebo} replaces the robot and the room: physics, simulated sensors
        (\pkg{/scan}, \pkg{/odom}, \pkg{/imu}, \pkg{/tf} odom$\rightarrow$base\_footprint), executes \pkg{/cmd\_vel}.
        Bridged to ROS via \pkg{ros\_gz\_bridge} (\pkg{turtlebot3\_gazebo/params/turtlebot3\_waffle\_pi\_bridge.yaml}).
  \item \textbf{Nav2} is the robot software: localization (AMCL), global planner, controller.
        It only sees what a real robot sees (laser, odometry), never the ground truth.
        In Task~2 Gazebo is replaced by the real TurtleBot, Nav2 stays the same.
  \item \textbf{RViz}: user interface and visualization only.
  \item \textbf{Ground truth}: known by Gazebo, but \emph{not} bridged to ROS by default
        (bridge: clock, joint\_states, odom, tf, cmd\_vel, imu, scan, camera).
        Needed only for evaluation (deviation metric). Solved in \pkg{mp2/launch/sim.launch.py}: PosePublisher
        plugin added to the robot model $\Rightarrow$ \pkg{/ground\_truth/pose} (30\,Hz).
  \item Map \pkg{turtlebot3\_navigation2/map/map.yaml}: 0.05\,m/cell, origin $(-10,-10)$;
        map frame origin = centre of the hexagonal \pkg{turtlebot3\_world}.
        Map frame and Gazebo world frame coincide: ground truth at the spawn pose $(-2.0006, -0.5000)$ equals the
        AMCL initial pose, and the localization error during the drives is 2--7\,cm without a systematic offset.
\end{itemize}

\subsection{Nav2 architecture}
\paragraph{TF chain and initial pose.} \pkg{map} $\rightarrow$ \pkg{odom} is published by \pkg{amcl},
\pkg{odom} $\rightarrow$ \pkg{base\_footprint} by the odometry (in simulation: Gazebo diff-drive plugin),
\pkg{base\_footprint} $\rightarrow$ \pkg{base\_link} by \pkg{robot\_state\_publisher}.
In \pkg{waffle\_pi.yaml}, \pkg{set\_initial\_pose: false}: AMCL waits for an initial pose
(RViz ``2D Pose Estimate'' or topic \pkg{/initialpose}) and publishes no \pkg{map} frame
until then. Until then the global costmap repeatedly logs
``Timed out waiting for transform from base\_link to map''. This is expected, not an error.
Robot spawn pose in \pkg{turtlebot3\_world}: $(x,y) = (-2.0, -0.5)$, yaw 0.

\paragraph{RViz.} RViz is only a client: it sends the initial pose (\pkg{/initialpose}) and the goal
(action \pkg{navigate\_to\_pose}) and visualizes topics. Nav2 runs without it.

\paragraph{Launch file chain.}
\begin{itemize}
  \item \pkg{turtlebot3\_navigation2/navigation2.launch.py}: only selects the files
        (map \pkg{map/map.yaml}, params \pkg{param/\$TURTLEBOT3\_MODEL.yaml} $=$ \pkg{waffle\_pi.yaml})
        and starts RViz. Everything else is included from \pkg{nav2\_bringup}.
  \item \pkg{nav2\_bringup/bringup\_launch.py}: \pkg{slam:=False} $\Rightarrow$ includes
        \pkg{localization\_launch.py} (otherwise \pkg{slam\_launch.py}) and \pkg{navigation\_launch.py}.
        \pkg{use\_composition:=True}: all nodes run as components in one process
        (\pkg{component\_container\_isolated}), hence the log prefix.
  \item \pkg{localization\_launch.py}: \pkg{map\_server}, \pkg{amcl} + own lifecycle manager.
  \item \pkg{navigation\_launch.py}: \pkg{controller\_server}, \pkg{smoother\_server}, \pkg{planner\_server},
        \pkg{route\_server}, \pkg{behavior\_server}, \pkg{velocity\_smoother}, \pkg{collision\_monitor},
        \pkg{bt\_navigator}, \pkg{waypoint\_follower}, \pkg{docking\_server} + own lifecycle manager.
\end{itemize}

\paragraph{Lifecycle nodes.} All Nav2 servers are lifecycle nodes
(unconfigured $\rightarrow$ inactive $\rightarrow$ active). The \pkg{lifecycle\_manager} configures and
activates them in the listed order, so that e.g.\ the costmaps only start when the map exists.

\paragraph{Nodes relevant for 1a.}
\begin{tabular}{p{3.2cm}p{5.2cm}p{7.2cm}}
\toprule
Node & Input $\rightarrow$ output & Role \\
\midrule
\pkg{map\_server} & \pkg{map.yaml} $\rightarrow$ \pkg{/map} & static occupancy grid \\
\pkg{amcl} & \pkg{/scan}, \pkg{/map}, tf odom $\rightarrow$ tf map$\rightarrow$odom & localization (MP1) \\
\pkg{bt\_navigator} & action \pkg{navigate\_to\_pose} & executes the behavior tree, calls planner and controller \\
\pkg{planner\_server} & action \pkg{compute\_path\_to\_pose} $\rightarrow$ \pkg{/plan} & global planner plugin (\pkg{GridBased} = NavFn) on the global costmap \\
\pkg{controller\_server} & action \pkg{follow\_path} $\rightarrow$ \pkg{cmd\_vel\_nav} & local planner plugin (\pkg{FollowPath} = DWB) on the local costmap, 10\,Hz \\
\pkg{velocity\_smoother} & \pkg{cmd\_vel\_nav} $\rightarrow$ \pkg{cmd\_vel\_smoothed} & limits velocity and acceleration, 20\,Hz \\
\pkg{collision\_monitor} & \pkg{cmd\_vel\_smoothed} $\rightarrow$ \pkg{cmd\_vel} & last safety stop/slowdown from raw laser data \\
\pkg{behavior\_server} & actions spin, backup, wait & recovery behaviors \\
\bottomrule
\end{tabular}
Global and local costmap run inside \pkg{planner\_server} and \pkg{controller\_server}.

\paragraph{Behavior tree} (\pkg{navigate\_to\_pose\_w\_replanning\_and\_recovery.xml}):
\begin{enumerate}
  \item \pkg{ComputePathToPose} at 1\,Hz (\pkg{RateController}) $\Rightarrow$ the global path is
        \emph{replanned every second} from the current pose, not only once.
  \item \pkg{FollowPath} runs in parallel (\pkg{PipelineSequence}) with the latest path.
  \item On failure: clear local/global costmap, then recoveries in round robin:
        spin $90^\circ$, wait 5\,s, back up 0.3\,m; up to 6 retries.
\end{enumerate}

\paragraph{Velocity command chain.}
controller (DWB) $\rightarrow$ \pkg{cmd\_vel\_nav} $\rightarrow$ velocity smoother $\rightarrow$
\pkg{cmd\_vel\_smoothed} $\rightarrow$ collision monitor $\rightarrow$ \pkg{cmd\_vel} $\rightarrow$ Gazebo / robot.

\paragraph{Key parameters (\pkg{waffle\_pi.yaml}).}
\begin{tabular}{lll}
\toprule
Parameter & Value & Meaning \\
\midrule
global costmap & static + obstacle + voxel + inflation, 1\,Hz & whole map \\
local costmap & obstacle + voxel + inflation, $3\times3$\,m rolling window, 5\,Hz & around the robot \\
\pkg{robot\_radius} & 0.15\,m & circular footprint \\
\pkg{inflation\_radius} / \pkg{cost\_scaling\_factor} & 0.5\,m / 5.0 & cost decay around obstacles \\
NavFn \pkg{tolerance} / \pkg{use\_astar} & 0.5\,m / false & goal tolerance of planner / Dijkstra \\
DWB \pkg{max\_vel\_x} / \pkg{max\_vel\_theta} & 0.3\,m/s / 1.0\,rad/s & \\
DWB \pkg{vx\_samples} / \pkg{vtheta\_samples} / \pkg{sim\_time} & 20 / 40 / 1.5\,s & sampled velocities, rollout horizon \\
goal checker \pkg{xy} / \pkg{yaw\_goal\_tolerance} & 0.25\,m / 0.25\,rad & when the goal counts as reached \\
\bottomrule
\end{tabular}

Note: \pkg{robot\_radius} 0.15\,m is smaller than the Waffle Pi footprint ($281\times306$\,mm $\Rightarrow$
circumscribed radius $\approx 0.21$\,m). Relevant with small inflation radius (part 1: 0.20\,m clearance would
touch a pillar with the real robot). \todo{Task 2: set \pkg{robot\_radius} $\approx 0.21$ or use a footprint polygon.}
Node graph: TikZ diagram in the progress presentation (\pkg{presentation\_mp2.tex}, slide ``Nav2 Architecture'').

\subsection{Global planner: Dijkstra vs.\ A*}
\textbf{Dijkstra (NavFn):} expands the grid from one end point in order of accumulated cost (distance + costmap
cost) until the other end point is reached. \textbf{A*} (\pkg{use\_astar: true}): same, but ordered by cost + a
distance heuristic to the other end point, so fewer cells should be expanded. \pkg{tolerance} (0.5\,m): if the goal
cell is blocked, the planner may end at a free cell within this distance.

\subsection{Local planner: DWB (DWA)}
\textbf{DWB / DWA:} samples velocity pairs $(v,\omega)$ within the dynamically reachable window
(\pkg{vx\_samples} 20, \pkg{vtheta\_samples} 40), simulates each for \pkg{sim\_time} 1.5\,s and picks the
trajectory with the best weighted score of the critics (PathDist, PathAlign, GoalDist, GoalAlign, BaseObstacle, ...).
\textbf{MPPI:} samples many noisy control sequences over a horizon, scores them with cost functions and takes their
cost-weighted average (installed alternative for the comparison).

\subsection{Costmaps}
\textbf{Costmap:} grid of costs 0--254 built from layers: static layer (map), obstacle/voxel layer (live laser),
inflation layer (cost decaying with distance to obstacles, $c(d) = 252\,e^{-\alpha(d - r_\text{robot})}$ for
$r_\text{robot} < d \le r_\text{infl}$). Global costmap: whole map, used by the planner; local costmap:
$3\times3$\,m rolling window around the robot, used by the controller.

\subsection{Planned experiments}
A script based on \pkg{nav2\_simple\_commander} drives the same set of start/goal pairs
for every configuration, so that the runs are reproducible and comparable. The same
script and metrics are reused in Task~1b for the RRT planner.

\paragraph{Metrics.} Few metrics that are easy to interpret (quality over quantity).
Planner metrics evaluate the global planner, the controller metric evaluates path following.
\begin{itemize}
  \item \textbf{Path length} $L = \sum_i \lVert \mathbf{p}_{i+1} - \mathbf{p}_i \rVert$ of the planned path [m] (planner)
  \item \textbf{Planning time} [ms] (planner); small differences expected for Dijkstra vs.\ A*
        in the small world, more relevant for RRT (random, several runs per goal needed)
  \item \textbf{Deviation} planned vs.\ driven path [m] (controller): for each driven pose the
        distance to the nearest point of the planned path; report mean and maximum.
        Computed with the estimated pose (controller error) and with the Gazebo ground truth
        (controller + localization error), see part 2.
  \item \textbf{Minimum clearance} to obstacles along the planned path [m] (planner).
        Path length alone rewards paths close to walls; clearance shows the trade-off
        short vs.\ safe that the inflation layer controls.
\end{itemize}
Not used: smoothness (visible in the plots), execution time (mostly follows from path
length), AMCL error (topic of MP1).

\paragraph{Which planned path is evaluated?} The default behavior tree replans every second,
so there is no single planned path. For RRT (1b) this would give a different random path every
second and the robot would jump between paths.
\textbf{Decision:} for all quantitative runs (1a and 1b) we use the behavior tree shipped with Nav2
\pkg{navigate\_w\_recovery\_and\_replanning\_only\_if\_path\_becomes\_invalid.xml}: plan once,
replan only if the goal changes or the path becomes blocked in the costmap (\pkg{IsPathValid}).
Same conditions for Dijkstra, A* and RRT; exactly one planned path per run, so path length,
clearance and deviation are well defined and the deviation measures only the controller.
The default replanning behavior (1\,Hz) is only explained/demonstrated in 1a.

\paragraph{Experiment design: dependence between global planner and controller.}
With the plan-once behavior tree, the global path is computed on the global costmap before the
robot moves, so the controller has no influence on it. Hence:
\begin{itemize}
  \item \textbf{Part 1 -- planner comparison without driving.} Paths are requested via
        \pkg{ComputePathToPose} with fixed start and goal (\pkg{use\_start: true}); the robot does not move.
        Metrics: path length, planning time, minimum clearance. Independent of the controller,
        fast $\Rightarrow$ many start/goal pairs; several repetitions per pair for RRT (random).
        \textbf{Runs without Gazebo} (\pkg{mp2/launch/planning\_only.launch.py}): only \pkg{map\_server},
        \pkg{planner\_server} and a static transform map$\rightarrow$base\_link. Reasons (found in a first test
        with the full simulation):
        \begin{enumerate}
          \item The global costmap also contains the obstacle and voxel layer, i.e.\ the \emph{current laser scan},
                placed with the AMCL pose. With a wrong AMCL pose, the same request gave a different path (P2: 3.32
                vs.\ 3.61\,m) or no path at all (P1, P3, error 208 \pkg{NO\_VALID\_PATH}). The planned path therefore
                depends on where the robot is and what it sees. Planning-only: static layer + inflation layer only
                $\Rightarrow$ identical costmap for every request.
          \item With \pkg{use\_sim\_time: true}, \pkg{planner\_server} measures the planning time with the
                Gazebo clock, which advances in steps (reported 0--5\,ms for the same request). Measuring in the
                client includes the action communication ($\approx$10\,ms without, $\approx$50\,ms with Gazebo running),
                which is larger than the planning itself. Planning-only: \pkg{use\_sim\_time: false}, the time
                reported by the server is wall-clock time of the planning alone.
        \end{enumerate}
        Same planner and costmap parameters as in the driving runs (the launch file reads \pkg{nav2\_params.yaml}).
  \item \textbf{Part 2 -- driving.} Reduced (decision 2026-11-10): \textbf{Dijkstra + DWB only}. A* gives
        almost the same paths, and comparing controllers is not required by the assignment. The controller
        is kept fixed at DWB; in 1b RRT is driven with DWB as well and compared with Dijkstra + DWB
        (assumption to state in the report: the deviation depends on planner \emph{and} controller).
        Setup (\pkg{mp2/launch/sim.launch.py}, \pkg{mp2/scripts/drive\_compare.py}):
        \begin{itemize}
          \item Gazebo headless, robot model \emph{without camera} (software rendering + raw image bridge
                overloaded the CPU under WSL2: real-time factor 0.3, scans every few seconds, Nav2 did not
                start), with a PosePublisher plugin $\Rightarrow$ ground truth \pkg{/ground\_truth/pose} at 30\,Hz.
          \item Per pair: drive to the start (not evaluated), then the measured run. Planned path = first
                \pkg{/plan} message of the run; estimated path = tf map$\rightarrow$base\_footprint at 10\,Hz;
                true path = ground truth.
          \item \textbf{Deviation (estimated)} = estimated pose vs.\ planned path: what the controller sees and
                corrects $\Rightarrow$ controller error. \textbf{Deviation (true)} = ground truth vs.\ planned
                path $\Rightarrow$ controller + localization error. Plus localization error (estimated vs.\ true).
                (Correction of an earlier statement: the ground truth does \emph{not} separate the controller
                error from the localization error, it contains both.)
          \item Goal checker: \pkg{yaw\_goal\_tolerance} 0.25 $\rightarrow$ 6.29\,rad (heading not checked).
                With 0.25, P3 failed: robot within 0.25\,m of the goal, rotates in place, DWB gets an empty
                pruned plan (``Resulting plan has 0 poses''), the plan-once tree does not replan (path still
                valid), recoveries, goal aborted. Weakness of plan-once; with 1\,Hz replanning this would not occur.
          \item \pkg{xy\_goal\_tolerance} 0.25\,m kept: start and end of a run differ from the pair coordinates
                by up to 0.25\,m, so the planned lengths differ from part 1.
        \end{itemize}
  \item \textbf{Other parameters} (\pkg{inflation\_radius}, \pkg{cost\_scaling\_factor}, \pkg{robot\_radius}):
        one at a time from the baseline (Dijkstra + DWB, default parameters), no full factorial design.
        Note: costmap parameters affect planner \emph{and} controller (both use a costmap).
\end{itemize}

\paragraph{Configurations (done).}
\begin{itemize}
  \item part 1: global planner NavFn Dijkstra vs.\ NavFn A*; costmap \pkg{inflation\_radius} 0.5/0.4/0.3/0.2\,m,
        \pkg{cost\_scaling\_factor} 2/5/10 (one at a time, Dijkstra)
  \item part 2: Dijkstra + DWB, default parameters (except \pkg{yaw\_goal\_tolerance}), 3 repetitions
\end{itemize}

\subsection{Results}
\paragraph{Part 1, baseline} (NavFn, default costmap: \pkg{robot\_radius} 0.15\,m, \pkg{inflation\_radius} 0.5\,m,
\pkg{cost\_scaling\_factor} 5.0; 20 requests per pair and planner; data \pkg{mp2/data/task1a/plans\_baseline.*},
figure \pkg{mp2/figures/plans\_baseline.png}).

\begin{tabular}{llccc}
\toprule
Pair & Planner & Length [m] & Planning time [ms] & Min.\ clearance [m] \\
\midrule
P1 (straight line 1.10\,m) & Dijkstra & 1.112 & $1.09 \pm 0.23$ & 0.350 \\
                           & A*       & 1.112 & $1.39 \pm 0.29$ & 0.350 \\
P2 (straight line 3.20\,m) & Dijkstra & 3.289 & $1.33 \pm 0.38$ & 0.320 \\
                           & A*       & 3.330 & $1.31 \pm 0.31$ & 0.320 \\
P3 (straight line 4.53\,m) & Dijkstra & 6.402 & $1.43 \pm 0.30$ & 0.335 \\
                           & A*       & 6.501 & $2.01 \pm 0.51$ & 0.335 \\
\bottomrule
\end{tabular}

Observations:
\begin{itemize}
  \item Both planners are deterministic: same path in all 20 requests (std of length 0).
  \item \textbf{P3: neither planner goes through the pillars}, both take the detour around the pillar field
        (6.4--6.5\,m instead of 4.5\,m), Dijkstra below/right, A* left/top.
        Hypothesis: NavFn minimizes the accumulated cost (distance + costmap cost). The gaps between the
        pillars are 0.8\,m wide, so every point in a gap is within 0.4\,m of a pillar, i.e.\ inside the inflation
        radius (0.5\,m) $\Rightarrow$ high cost. The outer ring is cheaper despite the longer distance.
        The two directions around are almost equally expensive (symmetric world), so small differences
        decide. \textbf{Confirmed} by the costmap sweep below.
  \item A* is \emph{not} faster than Dijkstra here (1--2\,ms for both). \todo{explain: small map, NavFn
        implementation (both stop when the start is reached?), overhead of the heuristic}
  \item Minimum clearance is the same for both planners (0.32--0.35\,m from the robot centre, i.e.\ 0.17--0.20\,m
        from the robot edge with \pkg{robot\_radius} 0.15\,m).
\end{itemize}

\paragraph{Part 1, costmap parameters} (Dijkstra, one parameter changed at a time, 10 requests each;
\pkg{mp2/scripts/run\_costmap\_sweep.sh}; figures \pkg{sweep\_inflation\_dijkstra.png}, \pkg{sweep\_csf\_dijkstra.png}).
Inflation cost for $r_\text{robot} < d \le r_\text{infl}$: $c(d) = 252\, e^{-\alpha (d - r_\text{robot})}$,
$\alpha$ = \pkg{cost\_scaling\_factor}; $c = 0$ for $d > r_\text{infl}$.
Planning time reported as median (single outliers from OS scheduling).

\begin{tabular}{lcccccc}
\toprule
 & \multicolumn{2}{c}{P1} & \multicolumn{2}{c}{P2} & \multicolumn{2}{c}{P3} \\
Configuration & $L$ [m] & clear.\ [m] & $L$ [m] & clear.\ [m] & $L$ [m] & clear.\ [m] \\
\midrule
$r_\text{infl}=0.5$, $\alpha=5$ (baseline) & 1.112 & 0.350 & 3.289 & 0.320 & 6.402 & 0.335 \\
$r_\text{infl}=0.4$ & 1.138 & 0.350 & 3.434 & 0.335 & 6.633 & 0.335 \\
$r_\text{infl}=0.3$ & 1.100 & 0.400 & 3.223 & 0.335 & 5.462 & 0.292 \\
$r_\text{infl}=0.2$ & 1.100 & 0.400 & 3.200 & 0.335 & 4.906 & 0.200 \\
$\alpha=2$  & 1.205 & 0.350 & 3.327 & 0.283 & 6.418 & 0.335 \\
$\alpha=10$ & 1.102 & 0.350 & 3.264 & 0.335 & 6.093 & 0.335 \\
\bottomrule
\end{tabular}

Observations:
\begin{itemize}
  \item \textbf{Hypothesis confirmed:} with $r_\text{infl} \le 0.3$\,m, the gaps between the pillars are no longer
        expensive and P3 goes through the pillar field (5.46\,m / 4.91\,m instead of 6.40\,m). Price: the clearance
        drops to 0.29 / 0.20\,m. 0.20\,m from the robot centre is only 0.05\,m from the edge with
        \pkg{robot\_radius} 0.15\,m -- and less than the real Waffle Pi radius ($\approx 0.21$\,m), so this path
        would touch the pillar with the real robot.
  \item $r_\text{infl}=0.4$ makes P3 and P2 \emph{longer} than the baseline (6.63 / 3.43\,m): the cost drops to 0
        abruptly at $d = r_\text{infl}$; the cheapest route changes in a non-monotonic way.
  \item P2 with small $r_\text{infl}$: almost the straight line (3.20\,m), the wall no longer pushes the path away.
  \item $\alpha$: larger $\alpha$ = cost decays faster = obstacles ``feel'' smaller: $\alpha=10$ shortens P3 to 6.09\,m
        (partly between the pillars), $\alpha=2$ makes P1 longer (1.21\,m) and P2 closer to the pillars (0.28\,m).
  \item Planning time 1--2\,ms in all configurations; differences are within the scatter.
        (Outlier $r_\text{infl}=0.4$/P3 with median 6\,ms: rerun gave 2.75\,ms with identical paths $\Rightarrow$
        measurement noise; table shows the rerun.) For NavFn the planning time is not a
        discriminating metric in this small map; it becomes relevant for RRT.
  \item Trade-off \textbf{short vs.\ safe}: inflation radius and cost scaling factor directly trade path length
        against clearance.
\end{itemize}

\paragraph{Part 2, driving} (Dijkstra + DWB, default parameters except \pkg{yaw\_goal\_tolerance}; 3 repetitions
of all pairs, 9/9 goals reached; data \pkg{mp2/data/task1a/drives\_dijkstra\_dwb\{,\_r2,\_r3\}.json},
figures \pkg{drives\_dijkstra\_dwb.png} (repetition 1), \pkg{drives\_dijkstra\_dwb\_all.png}).
Values: mean over the 3 repetitions of the per-run mean; max = largest value of all 3 runs.

\begin{tabular}{lcccccc}
\toprule
Pair & planned [m] & time [s] & dev.\ estimated [cm] & dev.\ true [cm] & localization err.\ [cm] \\
     &             &          & mean (max)            & mean (max)       & mean (max) \\
\midrule
P1 & 1.34 & 4.6  & 4.2 (11.4) & 2.9 (9.7)  & 2.2 (4.3) \\
P2 & 3.47 & 12.6 & 1.9 (5.8)  & 3.9 (10.7) & 5.0 (9.1) \\
P3 & 6.15 & 24.0 & 4.6 (22.9) & 5.8 (20.1) & 6.6 (11.7) \\
\bottomrule
\end{tabular}

Observations:
\begin{itemize}
  \item DWB follows the Dijkstra path within a few cm on average (1.9--4.6\,cm); the localization error (AMCL,
        2--7\,cm) is of the same order, so the true deviation is controller + localization error.
  \item Largest deviation in P3, always at the same place ($\approx(-1.5,-1.0)$, 4--5\,s after the start): the robot
        cuts $\approx$\,20\,cm to the right of the path (towards the bottom-left pillar) instead of following the path
        along the wall. Hypothesis: DWB critics trade path following (PathDist, PathAlign, scale 32) against progress
        to the goal (GoalDist, GoalAlign, scale 24). \todo{check, e.g.\ with higher PathDist.scale}
  \item P1: larger deviation in repetitions 2/3 (5--6\,cm) than in 1 (1\,cm): with \pkg{yaw\_goal\_tolerance} disabled,
        the robot reaches the start with an arbitrary heading and turns at the beginning of the run.
  \item Replanning: 1 plan per run, except repetition 1 / P2 (2 plans, path became invalid once).
  \item Robot stops 0.15--0.26\,m before the goal (\pkg{xy\_goal\_tolerance} 0.25\,m).
  \item First attempt (before fixing the recorder) showed 0.3--0.5\,m deviation in P2: partly caused by queued
        messages from the drive to the start (fixed by time-stamp filtering); not reproduced in 3 clean repetitions.
\end{itemize}

\subsection{Open questions (for the TA, progress meeting 2026-11-11)}
\begin{itemize}
  \item Is the planning-only setup (without Gazebo) acceptable for the planner comparison?
        Fallback: run the same requests once in the full simulation (robot localized at the start)
        and show that the paths are the same.
  \item Only DWB as controller (TEB not available in Jazzy) -- is a controller comparison expected?
  \item Why is A* not faster than Dijkstra in NavFn? \todo{}
\end{itemize}

\subsection{Pitfalls}
\begin{itemize}
  \item Leftover processes: \pkg{gz sim} can survive Ctrl+C and keeps publishing \pkg{/clock}, \pkg{/scan},
        \pkg{/tf}; leftover \pkg{static\_transform\_publisher} (map$\rightarrow$base\_link from the planning-only
        setup) split the TF tree of the simulation (``Timed out waiting for transform from base\_link to odom'').
        Check with \pkg{pgrep -fa "gz sim|ros2 launch|static\_transform"} before starting.
  \item Camera of the Waffle Pi: software rendering + raw image bridge overloaded the CPU under WSL2
        (real-time factor 0.3, Nav2 did not start) $\Rightarrow$ removed in \pkg{sim.launch.py}.
  \item Messages queue up while a node does not spin and are delivered later $\Rightarrow$ filter by time stamp.
  \item \pkg{nav2\_simple\_commander}: \pkg{BasicNavigator.waitUntilNav2Active()} publishes an initial pose at
        $(0,0)$ if it does not receive an \pkg{amcl\_pose}, which resets AMCL. Our scripts wait only for the
        node they need (\pkg{\_waitForNodeToActivate}).
\end{itemize}

\end{document}
